% kap09.tex — Klassisk feltteori
% Hamilton-Jacobi-kompendiet

\chapter{Klassisk feltteori}
\label{kap:felter}

Kompendiets sidste to kapitler forlader systemer med et \emph{endeligt}
antal frihedsgrader og bevæger sig mod \emph{felter} — systemer med et
\emph{uendeligt} (kontinuert) antal frihedsgrader, mærket af et
rumpunkt $x$ i stedet for et diskret indeks $i$. Dette er den nødvendige
forberedelse til kvantefeltteori, som er emnet for det senere,
selvstændige kompendium. Vi udvikler her Lagrange-billedet af klassisk
feltteori; Hamilton-billedet følger i kapitel~\ref{kap:felter-hamilton}.

\section{Fra diskret kæde til kontinuert felt}

\subsection{Den diskrete kæde}

Betragt $N$ masser, hver med masse $m$, forbundet i en lige kæde af
identiske fjedre med fjederkonstant $k$, i afstand $a$ fra hinanden i
ligevægt (se figur~\ref{fig:kaede-til-streng}, øverste del). Lad
$\varphi_i(t)$ betegne den $i$'te masses \emph{udsving} fra sin
ligevægtsposition (langs kæden, for simpelhedens skyld — en
"longitudinal" kæde). Lagrange-funktionen for dette system af koblede
oscillatorer er
\begin{equation}
  L = \sum_i\left[\frac12 m\dot\varphi_i^2 - \frac12 k(\varphi_{i+1}-\varphi_i)^2\right].
  \label{eq:L-diskret-kaede}
\end{equation}
Euler-Lagrange-ligningen for $\varphi_i$ (kapitel~\ref{kap:lagrange})
giver umiddelbart
\begin{equation}
  m\ddot\varphi_i = k(\varphi_{i+1}-\varphi_i) - k(\varphi_i-\varphi_{i-1}),
  \label{eq:kaede-bevaegelse}
\end{equation}
den velkendte ligning for en kæde af koblede oscillatorer (kraften på
masse $i$ fra dens to nabofjedre).

\subsection{Kontinuumsgrænsen}

Vi tager nu grænsen $a\to0$, $N\to\infty$, med den samlede længde
$L_{\text{kæde}}=Na$ fastholdt, og indfører de fysiske
\emph{tæthedsstørrelser}
\begin{equation}
  \mu \equiv \frac{m}{a} \quad(\text{masse pr.\ længde}), \qquad
  \tau \equiv ka \quad(\text{en effektiv "spænding"}),
\end{equation}
som vi holder fast, mens $a\to0$ (dvs.\ $m\to0$ og $k\to\infty$ på en
korreleret måde). Skriv \eqref{eq:L-diskret-kaede} om ved at
multiplicere og dividere med $a$ i hvert led:
\begin{equation}
  L = \sum_i a\left[\frac{1}{2}\frac{m}{a}\dot\varphi_i^2
      - \frac12 (ka)\left(\frac{\varphi_{i+1}-\varphi_i}{a}\right)^2\right]
    = \sum_i a\left[\frac12\mu\dot\varphi_i^2 - \frac12\tau\left(\frac{\Delta\varphi_i}{a}\right)^2\right].
\end{equation}
I grænsen $a\to0$ bliver $\varphi_i(t)\to\varphi(x,t)$ (en glat
funktion af den kontinuerte position $x$), $\Delta\varphi_i/a\to
\partial\varphi/\partial x$, og summen $\sum_i a(\cdots)$ bliver et
Riemann-integral $\int dx(\cdots)$:
\begin{equation}
  \boxed{L = \int dx\,\mathcal{L}, \qquad
  \mathcal{L}(\varphi,\dot\varphi,\varphi') = \frac12\mu\dot\varphi^2 - \frac12\tau(\varphi')^2,}
  \label{eq:L-tætthed-streng}
\end{equation}
hvor $\varphi'\equiv\partial\varphi/\partial x$. Funktionen
$\mathcal{L}$ kaldes \emph{Lagrange-tætheden}\index{Lagrange-tæthed}: den spiller for feltet
$\varphi(x,t)$ samme rolle, som $L(q,\dot q)$ spillede for en
partikel — men nu er "koordinaten" selve feltet, mærket af det
kontinuerte indeks $x$ i stedet for et diskret $i$.

\begin{figure}[htbp]
  \centering
  \begin{tikzpicture}[scale=1.0]
    % Diskret kæde
    \foreach \x/\lbl in {0/1,1.4/2,2.8/3,4.2/4,5.6/5} {
      \filldraw[fill=blue!20,draw=blue!60!black,thick] (\x,1.5) circle (0.18);
    }
    \foreach \x in {0,1.4,2.8,4.2} {
      \draw[decorate,decoration={coil,aspect=0.6,segment length=2.2mm,amplitude=1.8mm}]
        ($(\x,1.5)+(0.2,0)$) -- ($(\x+1.4,1.5)+(-0.2,0)$);
    }
    \node at (2.8,2.1) {diskret kæde ($a$, $m$, $k$)};
    % Pil ned
    \draw[-{Latex},thick] (2.8,1.0) -- (2.8,0.5);
    \node at (4.0,0.75) {$a\to0$};
    % Kontinuert streng
    \draw[thick,green!55!black] plot[smooth] coordinates
      {(0,-0.5) (0.9,-0.15) (1.8,-0.55) (2.8,-0.2) (3.8,-0.6) (4.8,-0.2) (5.6,-0.5)};
    \node[green!55!black] at (2.8,-1.1) {kontinuert streng ($\mu$, $\tau$)};
  \end{tikzpicture}
  \caption{Kontinuumsgrænsen: en diskret kæde af masser og fjedre
  (blå) bliver, når afstanden $a\to0$ med $\mu=m/a$ og $\tau=ka$
  fastholdt, til en kontinuert vibrerende streng (grøn), beskrevet af
  et felt $\varphi(x,t)$.}
  \label{fig:kaede-til-streng}
\end{figure}

\section{Hamiltons princip for felter}

Vi udleder nu feltets bevægelsesligning direkte fra et variations\-
princip, i fuld analogi med afsnit~\ref{sec:hamiltons-princip}. Lad
virkningen være
\begin{equation}
  S = \int_{t_1}^{t_2}dt\int dx\,\mathcal{L}(\varphi,\dot\varphi,\varphi',x,t),
\end{equation}
og betragt en variation $\varphi(x,t)\to\varphi(x,t)+\delta\varphi(x,t)$,
med $\delta\varphi=0$ ved $t=t_1,t_2$ og (for en endelig streng) ved
strengens endepunkter $x=0,L_{\text{kæde}}$. Variationen af $S$ er
\begin{equation}
  \delta S = \int dt\int dx\left(\pdv{\mathcal{L}}{\varphi}\delta\varphi
    + \pdv{\mathcal{L}}{\dot\varphi}\delta\dot\varphi
    + \pdv{\mathcal{L}}{\varphi'}\delta\varphi'\right).
\end{equation}
Vi integrerer de to sidste led partielt — med hensyn til $t$ for det
midterste led, og med hensyn til $x$ for det sidste — og bruger, at
$\delta\dot\varphi=\partial_t(\delta\varphi)$ og
$\delta\varphi'=\partial_x(\delta\varphi)$:
\begin{align}
  \int dt\,\pdv{\mathcal{L}}{\dot\varphi}\partial_t(\delta\varphi)
  &= \left[\pdv{\mathcal{L}}{\dot\varphi}\delta\varphi\right]_{t_1}^{t_2}
    - \int dt\,\dv{}{t}\left(\pdv{\mathcal{L}}{\dot\varphi}\right)\delta\varphi, \\
  \int dx\,\pdv{\mathcal{L}}{\varphi'}\partial_x(\delta\varphi)
  &= \left[\pdv{\mathcal{L}}{\varphi'}\delta\varphi\right]_{x=0}^{x=L_{\text{kæde}}}
    - \int dx\,\dv{}{x}\left(\pdv{\mathcal{L}}{\varphi'}\right)\delta\varphi.
\end{align}
Begge randtermer forsvinder — den første fordi $\delta\varphi=0$ ved
$t_1,t_2$, den anden fordi $\delta\varphi=0$ ved endepunkterne (eller,
for en uendelig streng, fordi $\varphi\to0$ tilstrækkeligt hurtigt for
$x\to\pm\infty$). Vi står tilbage med
\begin{equation}
  \delta S = \int dt\int dx\left[\pdv{\mathcal{L}}{\varphi}
    - \dv{}{t}\pdv{\mathcal{L}}{\dot\varphi} - \dv{}{x}\pdv{\mathcal{L}}{\varphi'}\right]\delta\varphi,
\end{equation}
og da $\delta\varphi(x,t)$ er arbitrær (variationsregningens
fundamentallemma, nu i to variable $x,t$ — samme argument som i
afsnit~\ref{sec:hamiltons-princip}, ledet igennem i hver variabel),
må koefficienten være identisk nul:
\begin{equation}
  \boxed{\dv{}{t}\pdv{\mathcal{L}}{\dot\varphi} + \dv{}{x}\pdv{\mathcal{L}}{\varphi'} - \pdv{\mathcal{L}}{\varphi} = 0.}
  \label{eq:euler-lagrange-felt}
\end{equation}
Dette er \emph{Euler-Lagrange-ligningen for et felt}\index{Euler-Lagrange-ligningerne!for felter} — den direkte
generalisering af \eqref{eq:euler-lagrange}, med et ekstra led for
hver rumlig retning, feltet afhænger af.

\section{Eksempel: bølgeligningen for strengen}

\begin{eksempel}
Vi anvender \eqref{eq:euler-lagrange-felt} på Lagrange-tætheden for
strengen, \eqref{eq:L-tætthed-streng}. De nødvendige afledte er
\begin{equation}
  \pdv{\mathcal{L}}{\dot\varphi} = \mu\dot\varphi, \qquad
  \pdv{\mathcal{L}}{\varphi'} = -\tau\varphi', \qquad
  \pdv{\mathcal{L}}{\varphi} = 0,
\end{equation}
så \eqref{eq:euler-lagrange-felt} giver
\begin{equation}
  \dv{}{t}(\mu\dot\varphi) + \dv{}{x}(-\tau\varphi') = 0
  \quad\Longrightarrow\quad
  \mu\ddot\varphi - \tau\varphi'' = 0,
\end{equation}
altså \emph{bølgeligningen}\index{bølgeligning}
\begin{equation}
  \boxed{\pdv[2]{\varphi}{t} = v^2\pdv[2]{\varphi}{x}, \qquad v\equiv\sqrt{\frac{\tau}{\mu}}.}
  \label{eq:boelgeligning}
\end{equation}
Vi har altså udledt bølgeligningen for en vibrerende streng —
inklusive den korrekte udbredelseshastighed $v=\sqrt{\tau/\mu}$, hvor
$\tau$ er strengens spænding og $\mu$ dens massefylde pr.\ længde —
helt uden at opskrive Newtons lov for et lille strengelement direkte.
Dette bekræfter, at Lagrange-tætheden \eqref{eq:L-tætthed-streng} er
den korrekte kontinuumsbeskrivelse af den diskrete kæde.
\end{eksempel}

\section{Noethers sætning for felter: energitæthed og kontinuitet}

Vi generaliserer nu Noether-resultatet fra
afsnit~\ref{eq:def-H-forlobig}'s diskussion (tidstranslationssymmetri
$\to$ bevaret $H$) til felter. Antag, at $\mathcal{L}$ ikke afhænger
eksplicit af $t$ (som for strengen). Vi beregner $d\mathcal{L}/dt$
langs en faktisk løsning, ved brug af kædereglen og derefter
Euler-Lagrange-ligningen \eqref{eq:euler-lagrange-felt} til at
omskrive $\partial\mathcal{L}/\partial\varphi$ (helt analogt til
udledningen, der gav \eqref{eq:def-H-forlobig}, men nu med et ekstra
led for den rumlige afledte):
\begin{align}
  \dv{\mathcal{L}}{t} &= \pdv{\mathcal{L}}{\varphi}\dot\varphi
    + \pdv{\mathcal{L}}{\dot\varphi}\ddot\varphi + \pdv{\mathcal{L}}{\varphi'}\dot\varphi' \\
  &= \left[\dv{}{t}\pdv{\mathcal{L}}{\dot\varphi} + \dv{}{x}\pdv{\mathcal{L}}{\varphi'}\right]\dot\varphi
    + \pdv{\mathcal{L}}{\dot\varphi}\ddot\varphi + \pdv{\mathcal{L}}{\varphi'}\dot\varphi' \\
  &= \dv{}{t}\left(\dot\varphi\pdv{\mathcal{L}}{\dot\varphi}\right)
    + \dv{}{x}\left(\pdv{\mathcal{L}}{\varphi'}\right)\dot\varphi + \pdv{\mathcal{L}}{\varphi'}\dv{}{x}(\dot\varphi) \\
  &= \dv{}{t}\left(\dot\varphi\pdv{\mathcal{L}}{\dot\varphi}\right)
    + \dv{}{x}\left(\dot\varphi\pdv{\mathcal{L}}{\varphi'}\right),
\end{align}
hvor vi i sidste skridt brugte $\dot\varphi'=\partial_x\dot\varphi$
(ombytning af blandede afledte, som i \eqref{eq:ombytning}) og
produktreglen baglæns endnu en gang. Omskrevet:
\begin{equation}
  \boxed{\pdv{\mathcal{H}}{t} + \pdv{\mathcal{S}}{x} = 0,}
  \qquad
  \mathcal{H} \equiv \dot\varphi\pdv{\mathcal{L}}{\dot\varphi} - \mathcal{L},
  \qquad
  \mathcal{S} \equiv \dot\varphi\pdv{\mathcal{L}}{\varphi'}.
  \label{eq:kontinuitetsligning}
\end{equation}
Dette er en \emph{kontinuitetsligning}\index{kontinuitetsligning} — den lokale form af en
bevarelseslov, velkendt fra elektromagnetismens ladningsbevarelse
($\partial\rho/\partial t+\partial J/\partial x=0$, i én rumlig
dimension) og fra hydrodynamikkens massebevarelse. $\mathcal{H}$
kaldes \emph{energitætheden}\index{energitæthed} og $\mathcal{S}$ \emph{energi\-
strømtætheden}\index{energistrømtæthed} (den effekt, der strømmer gennem et punkt $x$).

Integrerer vi \eqref{eq:kontinuitetsligning} over hele strengens
længde, og antager, at $\mathcal{S}\to0$ ved endepunkterne (eller ved
$x\to\pm\infty$ for en uendelig streng — fysisk: ingen effekt
strømmer ud af systemet), fås
\begin{equation}
  \dv{}{t}\int\mathcal{H}\,dx = -\Bigl[\mathcal{S}\Bigr]_{\text{rand}} = 0,
\end{equation}
altså at den \emph{totale} energi $E=\int\mathcal{H}\,dx$ er bevaret —
den kontinuerte analog til $H$'s bevarelse for et endeligt system
(kapitel~\ref{kap:hamilton}).

\subsection{Eksempel: energitæthed for strengen}

\begin{eksempel}
For strengen (med $\mathcal{L}$ fra \eqref{eq:L-tætthed-streng}) er
\begin{equation}
  \mathcal{H} = \dot\varphi(\mu\dot\varphi) - \left(\frac12\mu\dot\varphi^2-\frac12\tau(\varphi')^2\right)
  = \frac12\mu\dot\varphi^2 + \frac12\tau(\varphi')^2,
\end{equation}
den forventede sum af kinetisk energitæthed ($\frac12\mu\dot\varphi^2$)
og potentiel (elastisk) energitæthed ($\frac12\tau(\varphi')^2$, fra
strengens udspændthed). Energistrømtætheden er
\begin{equation}
  \mathcal{S} = \dot\varphi\cdot(-\tau\varphi') = -\tau\dot\varphi\,\varphi',
\end{equation}
som fysisk er den effekt, tværsnittet ved $x$ udfører på strengen til
højre for $x$ — en formel, der er velkendt fra strengens akustik og
fra bølgeledningsteori generelt.
\end{eksempel}

\section{Rumtranslationssymmetri: impulstæthed}
\label{sec:rum-translation}

Tidstranslation var kun det første af tre eksempler på Noethers
mønster "symmetri $\to$ bevaret strøm" for felter. Vi udleder nu det
andet: \emph{rumlig} translationssymmetri, som fører til
\emph{impulsbevarelse}, ved en udledning der er strukturelt identisk
med den foregående — blot med rollerne af $t$ og $x$ byttet om.

Antag stadig, at $\mathcal{L}$ ikke afhænger eksplicit af $x$ (som for
strengen: $\mu$ og $\tau$ er konstanter, uafhængige af position). Vi
beregner nu $d\mathcal{L}/dx$ langs en faktisk løsning, igen ved at
bruge Euler-Lagrange-ligningen \eqref{eq:euler-lagrange-felt} til at
omskrive $\partial\mathcal{L}/\partial\varphi$:
\begin{align}
  \dv{\mathcal{L}}{x} &= \pdv{\mathcal{L}}{\varphi}\varphi'
    + \pdv{\mathcal{L}}{\dot\varphi}\dot\varphi' + \pdv{\mathcal{L}}{\varphi'}\varphi'' \\
  &= \left[\dv{}{t}\pdv{\mathcal{L}}{\dot\varphi} + \dv{}{x}\pdv{\mathcal{L}}{\varphi'}\right]\varphi'
    + \pdv{\mathcal{L}}{\dot\varphi}\dot\varphi' + \pdv{\mathcal{L}}{\varphi'}\varphi'' \\
  &= \dv{}{t}\left(\varphi'\pdv{\mathcal{L}}{\dot\varphi}\right)
    + \dv{}{x}\left(\varphi'\pdv{\mathcal{L}}{\varphi'}\right),
\end{align}
hvor vi i sidste skridt igen brugte $\dot\varphi'=\partial_x\dot\varphi$
og produktreglen baglæns to gange, ganske som i det foregående afsnit
— blot med $\dot\varphi$ og $\varphi'$'s roller ombyttet. Omskrevet:
\begin{equation}
  \boxed{\pdv{\mathcal{P}}{t} + \pdv{\mathcal{T}}{x} = 0,}
  \qquad
  \mathcal{P} \equiv \varphi'\pdv{\mathcal{L}}{\dot\varphi},
  \qquad
  \mathcal{T} \equiv \varphi'\pdv{\mathcal{L}}{\varphi'} - \mathcal{L}.
  \label{eq:impuls-kontinuitet}
\end{equation}
$\mathcal{P}$ kaldes \emph{impulstætheden}\index{impulstæthed} og $\mathcal{T}$
\emph{impulsstrømtætheden}\index{impulsstrømtæthed} (eller \emph{spændingen}, i mekanisk
sprogbrug — det er netop den mekaniske spænding i materialet, betragtet
som en impulsstrøm). Integreret over hele strengen, med samme
antagelse om, at $\mathcal{T}\to0$ ved randen, giver dette bevarelse af
den \emph{totale} impuls $P=\int\mathcal{P}\,dx$.

\begin{bemaerkning}
Bemærk, at $\mathcal{P}=\varphi'\,\partial\mathcal{L}/\partial\dot\varphi
= \varphi'\pi$, hvor $\pi$ er den kanoniske impulstæthed, vi først
møder systematisk i kapitel~\ref{kap:felter-hamilton}. Dette er ikke
en tilfældighed: $\mathcal{P}$ og $\mathcal{H}$ (energitætheden) er de
to komponenter af det, der i relativistisk notation
(afsnit~\ref{sec:relativistisk-udblik}) kaldes
\emph{energi-impuls-tensoren}\index{energi-impuls-tensor} $T^{\mu\nu}$ — en enkelt, samlet
størrelse, hvis komponenter er netop de bevarede strømme, der svarer
til henholdsvis tids- og rumtranslation.
\end{bemaerkning}

\subsection{Eksempel: impulstæthed for strengen}

\begin{eksempel}
For strengen er
\begin{equation}
  \mathcal{P} = \varphi'(\mu\dot\varphi) = \mu\dot\varphi\,\varphi',
\end{equation}
og
\begin{equation}
  \mathcal{T} = \varphi'(-\tau\varphi') - \left(\frac12\mu\dot\varphi^2-\frac12\tau(\varphi')^2\right)
  = -\frac12\mu\dot\varphi^2 - \frac12\tau(\varphi')^2 = -\mathcal{H}.
\end{equation}
Det er en pæn, men ikke tilfældig detalje, at $\mathcal{T}=-\mathcal{H}$
for netop dette system: begge er specialtilfælde af den generelle
energi-impuls-tensor $T^{\mu\nu}=\partial\mathcal{L}/\partial(\partial_\mu\varphi)\,\partial^\nu\varphi-\eta^{\mu\nu}\mathcal{L}$
(afsnit~\ref{sec:relativistisk-udblik}), og for et fritfelt i $1+1$
dimensioner falder $T^{11}$ og $-T^{00}$ netop sammen på denne måde.
\end{eksempel}

\section{Intern symmetri: ladningsbevarelse}
\label{sec:ladning}

Det tredje og sidste eksempel på Noethers mønster er kvalitativt
forskelligt fra de to foregående: en symmetri, der \emph{ikke}
virker på rumtidskoordinaterne $t,x$ selv, men på feltets \emph{værdi}
— en såkaldt \emph{intern symmetri}\index{intern symmetri}. Dette kræver et felt med mindst
to reelle komponenter; det simplest mulige eksempel er et
\emph{komplekst} felt $\psi(x,t)$ (fysisk: to reelle felter, pakket
sammen som real- og imaginærdel).

Betragt Lagrange-tætheden
\begin{equation}
  \mathcal{L} = \mu\,\dot\psi^\ast\dot\psi - \tau\,\psi^{\ast\prime}\psi',
  \label{eq:L-komplekst-felt}
\end{equation}
en naturlig komplex generalisering af strengens Lagrange-tæthed
\eqref{eq:L-tætthed-streng}, hvor $\psi$ og $\psi^\ast$ nu behandles
som to \emph{uafhængige} felter i variationsprincippet (dette er den
sædvanlige og fuldt korrekte måde at behandle et komplekst felt — se
opgave 3). Euler-Lagrange-ligningerne for $\psi$ og $\psi^\ast$ giver
hver sin bølgeligning, af samme form som \eqref{eq:boelgeligning}.

$\mathcal{L}$ i \eqref{eq:L-komplekst-felt} er uændret under den
\emph{globale fasetransformation}\index{fasetransformation, global}
\begin{equation}
  \psi \to e^{i\alpha}\psi, \qquad \psi^\ast \to e^{-i\alpha}\psi^\ast,
  \qquad \alpha\in\mathbb{R} \text{ konstant},
\end{equation}
fordi hvert led i $\mathcal{L}$ indeholder netop ét $\psi$ og ét
$\psi^\ast$, så faktorerne $e^{i\alpha}$ og $e^{-i\alpha}$ altid
ophæver hinanden. Denne symmetri kaldes en \emph{$U(1)$-symmetri}\index{U(1)-symmetri}. For en
infinitesimal transformation, $\alpha=\varepsilon\ll1$:
\begin{equation}
  \delta\psi = i\varepsilon\psi, \qquad \delta\psi^\ast = -i\varepsilon\psi^\ast.
\end{equation}

\subsection{Udledning af den bevarede ladningsstrøm}

Invarians af $\mathcal{L}$ betyder, at $\delta\mathcal{L}=0$ til første
orden i $\varepsilon$. Vi udregner $\delta\mathcal{L}$ ved kædereglen
og bruger Euler-Lagrange-ligningerne for $\psi$ og $\psi^\ast$ til at
omskrive $\partial\mathcal{L}/\partial\psi$ og
$\partial\mathcal{L}/\partial\psi^\ast$ — helt analogt til de to
foregående udledninger, men nu med \emph{to} felter i stedet for et:
\begin{align}
  0 = \delta\mathcal{L} &= \pdv{\mathcal{L}}{\psi}\delta\psi
    + \pdv{\mathcal{L}}{\dot\psi}\delta\dot\psi + \pdv{\mathcal{L}}{\psi'}\delta\psi'
    + \pdv{\mathcal{L}}{\psi^\ast}\delta\psi^\ast
    + \pdv{\mathcal{L}}{\dot\psi^\ast}\delta\dot\psi^\ast
    + \pdv{\mathcal{L}}{\psi^{\ast\prime}}\delta\psi^{\ast\prime} \\
  &= \dv{}{t}\left(\pdv{\mathcal{L}}{\dot\psi}\delta\psi + \pdv{\mathcal{L}}{\dot\psi^\ast}\delta\psi^\ast\right)
   + \dv{}{x}\left(\pdv{\mathcal{L}}{\psi'}\delta\psi + \pdv{\mathcal{L}}{\psi^{\ast\prime}}\delta\psi^\ast\right),
\end{align}
hvor mellemskridtet (samme type omgruppering som i de to foregående
afsnit, anvendt på hvert af de fire led med $\delta\psi'$/$\delta\dot\psi$
osv.) er udført i detaljer i opgave 3. Indsættes
$\delta\psi=i\varepsilon\psi$, $\delta\psi^\ast=-i\varepsilon\psi^\ast$,
og divideres med $\varepsilon$, fås
\begin{equation}
  \boxed{\pdv{\rho}{t} + \pdv{j}{x} = 0,}
  \qquad
  \rho \equiv i\left(\psi\pdv{\mathcal{L}}{\dot\psi} - \psi^\ast\pdv{\mathcal{L}}{\dot\psi^\ast}\right),
  \qquad
  j \equiv i\left(\psi\pdv{\mathcal{L}}{\psi'} - \psi^\ast\pdv{\mathcal{L}}{\psi^{\ast\prime}}\right).
  \label{eq:ladnings-kontinuitet}
\end{equation}
$\rho$ kaldes \emph{ladningstætheden}\index{ladningstæthed} og $j$ \emph{ladningsstrøm\-
tætheden}\index{ladningsstrømtæthed}. Integreret over hele rummet (med $j\to0$ ved randen) giver
dette bevarelse af den \emph{totale ladning} $Q=\int\rho\,dx$ —
en tredje, kvalitativt ny type bevarelseslov, der ikke stammer fra en
rumtidssymmetri, men fra feltets egen interne struktur.

\subsection{Eksempel: ladningstæthed for det komplekse felt}

\begin{eksempel}
For \eqref{eq:L-komplekst-felt} er
$\partial\mathcal{L}/\partial\dot\psi=\mu\dot\psi^\ast$ og
$\partial\mathcal{L}/\partial\dot\psi^\ast=\mu\dot\psi$, så
\begin{equation}
  \rho = i\mu\left(\psi\dot\psi^\ast - \psi^\ast\dot\psi\right) = -2\mu\,\text{Im}(\psi^\ast\dot\psi),
\end{equation}
hvor vi brugte, at $\psi\dot\psi^\ast-\psi^\ast\dot\psi=-2i\,\text{Im}(\psi^\ast\dot\psi)$
for et komplekst tal. Tilsvarende er ladningsstrømtætheden
\begin{equation}
  j = i\left(\psi(-\tau\psi^{\ast\prime}) - \psi^\ast(-\tau\psi')\right)
    = -i\tau\left(\psi\psi^{\ast\prime}-\psi^\ast\psi'\right) = 2\tau\,\text{Im}(\psi^\ast\psi').
\end{equation}

\begin{bemaerkning}
Formen $\rho\propto\text{Im}(\psi^\ast\dot\psi)$ bør give genklang: det
er strukturelt \emph{samme} udtryk som sandsynlighedstætheden og
-strømmen, man møder for Schrödinger-ligningens bølgefunktion i
[[kvantemekanik-kompendiet]] ($\rho_{\text{QM}}\propto\text{Im}(\psi^\ast\dot\psi)$
eller, i den sædvanlige rumlige form,
$j_{\text{QM}}\propto\text{Im}(\psi^\ast\nabla\psi)$). Dette er ingen
tilfældighed: begge er Noether-ladningen for netop den samme $U(1)$-
fasesymmetri, blot anvendt på to forskellige felter — et klassisk
mekanisk felt her, og kvantemekanikkens bølgefunktion der. Sammenhængen
er et af de tydeligste eksempler i hele kompendiet på, hvordan
klassisk feltteoris begrebsapparat går næsten upåvirket over i
kvantemekanikken.
\end{bemaerkning}
\end{eksempel}

\begin{figure}[htbp]
  \centering
  \begin{tikzpicture}[scale=1.0]
    \draw[-{Latex}] (-0.3,0) -- (7.0,0) node[right] {$x$};
    % Kontrolvolumen
    \draw[thick,blue!60!black,fill=blue!8] (2.2,-0.9) rectangle (4.8,0.9);
    \node[blue!60!black] at (3.5,1.3) {kontrolvolumen $[x_1,x_2]$};
    % Ind- og udstrømning
    \draw[-{Latex},thick,green!55!black] (0.6,0) -- (2.0,0);
    \node[green!55!black] at (1.3,0.35) {ind};
    \draw[-{Latex},thick,green!55!black] (5.0,0) -- (6.4,0);
    \node[green!55!black] at (5.7,0.35) {ud};
    % Densitet inde i boksen (symbolsk)
    \node at (3.5,0) {$\displaystyle\int_{x_1}^{x_2}\!\rho\,dx$};
  \end{tikzpicture}
  \caption{Det fælles mønster bag alle tre bevarelseslove i dette
  kapitel: en kontinuitetsligning $\partial\rho/\partial t+\partial j/\partial x=0$
  siger, at ændringen i den samlede mængde ($\rho$: energi, impuls
  eller ladning) inde i et område kun kan skyldes strøm ($j$) ind eller
  ud gennem randen — intet forsvinder eller opstår "midt i" volumenet.}
  \label{fig:kontinuitet-generel}
\end{figure}

\section{Udblik: relativistisk notation og Klein-Gordon-feltet}
\label{sec:relativistisk-udblik}

Vi antyder her, hvordan formalismen generaliserer til relativistiske
felter i tre rumlige dimensioner — det udgangspunkt, det senere
QFT-kompendium starter fra. Med rumtidskoordinater
$x^\mu=(ct,x,y,z)$, $\mu=0,1,2,3$, og notationen
$\partial_\mu\varphi\equiv\partial\varphi/\partial x^\mu$, skrives
Lagrange-tætheden som $\mathcal{L}(\varphi,\partial_\mu\varphi)$, og
Euler-Lagrange-ligningen \eqref{eq:euler-lagrange-felt} generaliserer
kompakt til
\begin{equation}
  \partial_\mu\left(\pdv{\mathcal{L}}{(\partial_\mu\varphi)}\right) - \pdv{\mathcal{L}}{\varphi} = 0
  \qquad\text{(summeret over } \mu = 0,1,2,3\text{).}
\end{equation}
Det kanoniske eksempel på et relativistisk felt er
\emph{Klein-Gordon-Lagrangianen}\index{Klein-Gordon-Lagrangianen},
\begin{equation}
  \mathcal{L}_{\text{KG}} = \frac12\partial_\mu\varphi\,\partial^\mu\varphi - \frac12 m^2\varphi^2,
\end{equation}
som giver \emph{Klein-Gordon-ligningen}\index{Klein-Gordon-ligningen} $(\Box+m^2)\varphi=0$ (med
$\Box\equiv\partial^\mu\partial_\mu$) — den relativistiske
bølgeligning for et skalarfelt med masse $m$. Bemærk den strukturelle
lighed med strengens Lagrange-tæthed \eqref{eq:L-tætthed-streng}: begge
er af formen "kinetisk led minus gradientled (minus et eventuelt
masseled)". Denne parallel er ikke overfladisk — det er præcis
udgangspunktet for at kvantisere feltet $\varphi$ (kanonisk
feltkvantisering), hvilket er det senere QFT-kompendiums fundament.

\begin{bemaerkning}
Overgangen fra den mekaniske streng til Klein-Gordon-feltet er et
eksempel på et generelt mønster i fysikken: den samme matematiske
struktur (en Lagrange-tæthed, kvadratisk i felt og afledte)
beskriver både lydbølger i et fast stof \emph{og} kvantefelter i
vakuum. Forskellen ligger ikke i formalismen, men i, hvad "feltet"
$\varphi$ fysisk repræsenterer, og i hvordan det efterfølgende
kvantiseres. Vi vender ikke tilbage til dette spor i dette kompendium,
men det er værd at have set strukturen én gang i sin simpleste,
mekaniske form, før man møder den igen i en relativistisk og
kvantiseret kontekst.
\end{bemaerkning}

\section{Opsummering}

\begin{itemize}
  \item En diskret kæde af koblede oscillatorer bliver i
        kontinuumsgrænsen ($a\to0$, $\mu=m/a$, $\tau=ka$ fastholdt)
        til et felt $\varphi(x,t)$ med Lagrange-tæthed
        $\mathcal{L}=\frac12\mu\dot\varphi^2-\frac12\tau(\varphi')^2$.
  \item Hamiltons princip for felter giver Euler-Lagrange-ligningen
        $\dv*{}{t}(\partial\mathcal{L}/\partial\dot\varphi) +
        \dv*{}{x}(\partial\mathcal{L}/\partial\varphi') -
        \partial\mathcal{L}/\partial\varphi = 0$.
  \item For strengen giver dette bølgeligningen
        $\ddot\varphi=v^2\varphi''$, $v=\sqrt{\tau/\mu}$.
  \item Tidstranslationssymmetri giver en kontinuitetsligning
        $\partial\mathcal{H}/\partial t + \partial\mathcal{S}/\partial x=0$,
        med energitæthed $\mathcal{H}=\dot\varphi\,\partial\mathcal{L}/\partial\dot\varphi-\mathcal{L}$
        og energistrøm $\mathcal{S}=\dot\varphi\,\partial\mathcal{L}/\partial\varphi'$.
  \item Rumtranslationssymmetri giver analogt
        $\partial\mathcal{P}/\partial t+\partial\mathcal{T}/\partial x=0$,
        med impulstæthed $\mathcal{P}=\varphi'\,\partial\mathcal{L}/\partial\dot\varphi$
        og impulsstrøm (spænding)
        $\mathcal{T}=\varphi'\,\partial\mathcal{L}/\partial\varphi'-\mathcal{L}$.
        $\mathcal{P}$ og $\mathcal{H}$ er komponenter af
        energi-impuls-tensoren $T^{\mu\nu}$.
  \item En intern $U(1)$-fasesymmetri af et komplekst felt $\psi$
        giver en tredje kontinuitetsligning,
        $\partial\rho/\partial t+\partial j/\partial x=0$, med
        ladningstæthed $\rho=i(\psi\,\partial\mathcal{L}/\partial\dot\psi-\psi^\ast\,\partial\mathcal{L}/\partial\dot\psi^\ast)$
        — strukturelt identisk med kvantemekanikkens
        sandsynlighedstæthed.
  \item Alle tre bevarelseslove (energi, impuls, ladning) er
        specialtilfælde af samme mønster: en lokal kontinuitetsligning
        $\partial\rho/\partial t+\partial j/\partial x=0$, der
        integreret giver en global bevaret størrelse.
  \item Relativistisk generaliseres formalismen via $\partial_\mu$;
        Klein-Gordon-Lagrangianen har samme struktur som strengens,
        og peger direkte mod det senere QFT-kompendium.
\end{itemize}

\begin{opgave}
Udfør kontinuumsgrænsen for den \emph{transversale} kæde (masser
forbundet i en lige linje, men med udsving $\varphi_i$ vinkelret på
kædens retning, under en fast spænding $\tau_0$ mellem hver
nabomasse), og vis, at man igen når frem til
\eqref{eq:L-tætthed-streng} med $\tau=\tau_0$ (denne gang er $\tau$
den fysiske spænding direkte, ikke $ka$) — dette er den mere
sædvanlige model for en vibrerende (guitar-)streng.
\end{opgave}

\begin{opgave}
Vis, ved direkte indsættelse i den relativistiske Euler-Lagrange-
ligning, at Klein-Gordon-Lagrangianen $\mathcal{L}_{\text{KG}}$
faktisk giver $(\Box+m^2)\varphi=0$. (Brug metrikken
$\eta^{\mu\nu}=\text{diag}(1,-1,-1,-1)$, så
$\partial_\mu\varphi\partial^\mu\varphi = (\partial_t\varphi)^2/c^2 - |\nabla\varphi|^2$.)
\end{opgave}

\begin{opgave}
Udfør mellemskridtet i udledningen af \eqref{eq:ladnings-kontinuitet}
i detaljer: brug Euler-Lagrange-ligningerne for $\psi$ og $\psi^\ast$
til at omskrive $\partial\mathcal{L}/\partial\psi\cdot\delta\psi$ og
$\partial\mathcal{L}/\partial\psi^\ast\cdot\delta\psi^\ast$, og vis, at
de resterende led kan grupperes til en total tids- og en total
rumafledt, ganske som i udledningerne af
\eqref{eq:kontinuitetsligning} og \eqref{eq:impuls-kontinuitet}.
\end{opgave}

\begin{opgave}
Skriv $\psi=\varphi_1+i\varphi_2$ med $\varphi_1,\varphi_2$ reelle
felter, og vis, at Lagrange-tætheden \eqref{eq:L-komplekst-felt}
bliver $\mathcal{L}=\mu(\dot\varphi_1^2+\dot\varphi_2^2)-\tau(\varphi_1'^2+\varphi_2'^2)$
— to uafhængige, reelle strenge. Vis derefter, at $U(1)$-transformationen
$\psi\to e^{i\alpha}\psi$ svarer til en \emph{rotation} af
$(\varphi_1,\varphi_2)$ i "internt rum", og at den bevarede ladning
$Q=\int\rho\,dx$ er den kontinuerte analog til banemomentet
$L_z=xp_y-yp_x$ fra kapitel~\ref{kap:poisson} — nu med $(\varphi_1,\varphi_2)$
i stedet for $(x,y)$.
\end{opgave}
